\documentclass[11pt,a4paper]{article}
\usepackage[margin=19mm,headheight=15pt]{geometry}
\usepackage{fontspec}
\setmainfont{Noto Serif}
\usepackage{amsmath,amssymb,mathtools}
\usepackage{xcolor}
\usepackage{fancyhdr}
\usepackage{enumitem}
\usepackage{microtype}
\setlength{\parindent}{0pt}
\setlength{\parskip}{6pt}
\definecolor{navy}{HTML}{183153}
\definecolor{lightgray}{HTML}{F3F5F7}
\pagestyle{fancy}\fancyhf{}
\lhead{\small ME 5309 -- Mid-Semester Examination}
\rhead{\small ME 5309}
\cfoot{\thepage}
\begin{document}
\begin{center}
{\Large\bfseries ME 5309 -- Mid-Semester Examination}\\[3pt]
{\large Answer Booklet}\\[8pt]
\end{center}
\textbf{Instructions:} Attempt any three questions. Each question carries 10 marks. Show working where appropriate.
\section*{Question 1 \hfill [10 marks]}
\subsection*{Q1(a) \hfill [4 marks]}

The pixel pitch is 0.005 mm, so f\_x = 16/0.005 = 3200 pixels.
Sensor width = 6.4 mm.
Belt field of view = 400x6.4/16 = 160 mm; spatial resolusion = 160/1280 = 0.125 mm/pixel.

\subsection*{Q1(b) \hfill [3 marks]}

ADC step = 5/4096 = 1.2207 mV.
I divide by 0.1 mV/N, giving ΔF = 12.207 N.

\subsection*{Q1(c) \hfill [3 marks]}

\vspace{2mm}

\section*{Question 3 \hfill [10 marks]}
\subsection*{Q3(a) \hfill [2 marks]}

The estimate is noisy and can fluctuate from sample to sample.

\subsection*{Q3(b) \hfill [2 marks]}

For H(s)=ωc/(s+ωc), ωc=2π(20)=125.664 rad/s and T=0.001 s.
Tustin gives H(z)=α(1+z⁻¹)/[(1+α)+(α-1)z⁻¹], α=ωcT/2=0.0628319.
Thus:\[\begin{aligned}y[k]=0.881765y[k-1]+0.0591175(x_k+x[k-1])\end{aligned}\]
DC gain = (2x0.0591175)/(1-0.881765) approximately 1.000.\par\smallskip $\omega_c=2\pi(20)=125.664$ rad/s. With $T=0.001$ s, $\alpha=\omega_cT/2=125.664(0.001)/2=0.062832$. The recurrence coefficients are approximately $0.881765$ on the previous output and $0.0591175$ on each input sample. Check: $2(0.0591175)/(1-0.881765)\approx1$, so the DC gain is close to one.

\subsection*{Q3(c) \hfill [4 marks]}

A-LC=[[-ℓ₁,1], [-ℓ₂,-2]].
I obtain p(s)=s²+(ℓ₁+2)s+(ℓ₁+ℓ₂).
Matching (s+10)² gives ℓ₁=18 and ℓ₂=82.

\subsection*{Q3(d) \hfill [2 marks]}

The observer is simpler because it only smooths the measured velocity.
The filter is safer when the model is poor.

\section*{Question 4 \hfill [10 marks]}
\subsection*{Q4(a) \hfill [2 marks]}

The characteristic polynomial is s²+4s+8, so I report poles -4 and -8.
Unity DC gain requires Nff=8.

\subsection*{Q4(b) \hfill [2 marks]}

\vspace{2mm}

\subsection*{Q4(c) \hfill [4 marks]}

The desired polynomial is (s+4)(s²+4s+8).
The closed-loop polynomial should be cubic because an integrator is added.

\subsection*{Q4(d) \hfill [2 marks]}

Integral action helps eliminate error.
A problem is that the controller can become unstable.

\end{document}